\begin{theorem}[G-26: The complete first-row padding boundary through $m$ rows]\label{G-26}
Let $a=n-m>0$. For every coordinate degree $d\ge1$ and partition
$\lambda\vdash nd$ of length at most $r$, the padded permanent
multiplicity satisfies
\begin{equation}\label{G26:pieri_bound}
 \begin{cases}
 m_\lambda(\ell^a\operatorname{per}_m)=0,&\lambda_1<ad,\\[2pt]
 m_\lambda(\ell^a\operatorname{per}_m)
   \le m_{(\lambda_2,\ldots,\lambda_r)}(\operatorname{per}_m),
       &\lambda_1=ad.
 \end{cases}
\end{equation}
The multiplicity on the second line can be computed on the
permanent image in $r-1$ variables. This bound holds in every
finite degree over $\mathbb C$; no lifting or equality of these
multiplicity spaces is asserted.

Now suppose $n>m\ge17$ and $4\le r\le m$. Fix
$\nu\vdash ns$, $s\ge1$, with exactly $r$ positive parts.
If $\nu_1=(n-m)s$, the determinant multiplicity is strictly
larger than the padded permanent multiplicity at every
sufficiently large integer stretch $j$. The same conclusion
holds when $\nu_1<(n-m)s$, in which case the padded multiplicity
is zero at every positive stretch.

Consequently, in this row range every fixed full-row direction
has an exact eventual comparison: the padded permanent
multiplicity is eventually larger if and only if
\[
 \nu_1>(n-m)s
 \quad\hbox{and}\quad
 rm^2-(r-2)n^2-2m+r-1>0.
\]
In every other case the determinant multiplicity is eventually
strictly larger. In the strict-padding case the bounds are
those of G-23; an explicit equality-face bound follows below.
\end{theorem}

\paragraph{An explicit threshold on the equality face.}
Assume $\nu_1=as$, put $q=r-1$, and define
\[
 \begin{aligned}
 \mu&=(\nu_2,\ldots,\nu_r),&
 \eta&=\nu-\nu_r\mathbf1_r,&
 \theta&=\mu-\nu_r\mathbf1_q,\\
 e&=|\eta|,&e'&=|\theta|,&
 r_\eta&=\#\{i<j:\eta_i>\eta_j\},
 \end{aligned}
\]
with $r_\theta$ defined in the same way. Set
\[
 \begin{aligned}
 u&=(r-2)n^2-r^2+3+r_\eta,\\
 v&=\begin{cases}
       \binom{m+2}{2}-8+r_\theta,&r=4,\\
       (r-1)m^2-2m-r^2+2r+2+r_\theta,&r\ge5,
      \end{cases}\\
 b_D&=\binom{\binom{n+r-1}{r-1}+s-1}{s}\dim S_\eta\mathbb C^r,\\
 b_B&=\binom{\binom{m+q-1}{q-1}+s-1}{s}\dim S_\theta\mathbb C^q,\\
 b&=\max\{b_D,b_B\},&
 L&=ns+(r-1)e,&p&=ns+e,\\
 \beta&=\left\lceil\frac{3r^2b}{8}\right\rceil L^{2(r^2-1)},&
 T&=\binom{b+\beta}{b}-1,& E&=rT\beta!.
 \end{aligned}
\]
Then $u>v\ge2$, and strict determinant dominance holds for every
\begin{equation}\label{G26:threshold}
                  j\ge E\bigl(4v(pE)^{v-1}+1\bigr),
\end{equation}
in original coordinate degree $sj$. All Schur dimensions here
are given by the Weyl products used in G-23.

\paragraph{Background used.}
The Pieri formula is cited from Bump and Gamburd (2006),
\S2.1, equation with source label \texttt{e:pieri}, supplied
\texttt{acp20.tex}, lines 404--422. That source identifies it
with Macdonald's equation I~(5.16), p.~73, and Bump's
Theorem~42.4; we inspected the supplied Bump--Gamburd statement.
All other inputs are the exact invariant-extraction, degree,
and Hilbert bounds already used and located in G-22--G-24.

\begin{proof}
We first prove the finite-degree inequality and its change of
ambient dimension. We then compare two different enhanced
invariant algebras, with their degree semigroups treated separately.

\paragraph{1. Pieri removes exactly the first row.}
Let $V=\mathbb C^r$ and let $Y$ be the padded permanent image.
The parameter morphism
\[
 V\times\mathcal P_{m,r}\longrightarrow Y,
                    \qquad(\ell,F)\longmapsto\ell^a F
\]
has dense image by the definition of $Y$. Pullback of degree-$d$
coordinate functions is therefore an injective equivariant map
\begin{equation}\label{G26:parameter_injection}
 \mathbb C[Y]_d\lhook\joinrel\longrightarrow
       \operatorname{Sym}^{ad}V^*\otimes\mathbb C[\mathcal P_{m,r}]_d.
\end{equation}
The two parameter degrees are $(ad,d)$; they are not both $d$.
Use the contragredient convention for coordinate representations
fixed in the introduction. By complete reducibility over
$\mathbb C$ and Pieri, a copy of $S_\lambda V^*$ in the right
side comes from $S_\alpha V^*$ precisely when $\lambda/\alpha$
is a horizontal strip of size $ad$, with multiplicity one.

There are only $\lambda_1$ columns. If $\lambda_1<ad$, no such
strip exists. If $\lambda_1=ad$, the strip must contain exactly
one box in every column. Thus
\[
 \alpha'_k=\lambda'_k-1\quad(1\le k\le\lambda_1),
 \qquad
                 \alpha=(\lambda_2,\ldots,\lambda_r).
\]
This is the unique possibility, and its size is $md$.
Taking multiplicities in \eqref{G26:parameter_injection} gives
the asserted zero and upper bound.

For completeness, the multiplicity of this tail partition
equals that on $\mathcal P_{m,r-1}$. A highest-weight coordinate
polynomial with last weight zero uses only coefficient variables
of monomials not involving the last original variable; all these
coefficient weights are nonnegative in the polynomial convention.
Its vanishing on the $r$-variable permanent image is equivalent
to vanishing on the $(r-1)$-variable image: restriction of every
$r$-variable matrix pencil gives a pencil with one fewer variable,
and every such smaller pencil occurs by setting the unused
coefficient matrix to zero. The highest-weight coefficient
subalgebras and their ideal intersections consequently agree.
This is precisely the inheritance argument stated in the
introduction, now applied to the tail shape of length at most
$r-1$. It proves the last assertion of
\eqref{G26:pieri_bound} without replacing an actual coordinate
ring by its ambient plethysm ring.

\paragraph{2. The equality face gives two stable comparison algebras.}
Now assume the displayed equality-face hypotheses. Because
$\nu_1=as$ is the largest of the $r$ parts,
\[
 ns\le ras,\qquad
 n\ge\frac{rm}{r-1},\qquad
 a\ge\frac{m}{r-1}>1.
\]
Thus the integer $a$ is at least two. The tail $\mu$ has size
$ms$, exactly $q=r-1$ positive parts, and
\[
                    \mu_1\ge\frac{ms}{r-1}>s.
\]
Let $A$ be the enhanced invariant coordinate algebra for
the determinant cone $\mathcal D_{n,r}$, outer stretch $s$,
and flag weight $\eta$, as in G-22. Then
\[
                       \dim A_j=m_{j\nu}(\det_n).
\]
For $r\ge5$, let $B$ be the corresponding enhanced invariant
algebra for $\mathcal P_{m,q}$ and flag weight $\theta$.
For $r=4$, take instead the enhanced algebra for the full
ternary degree-$m$ coefficient space. In that case exactness
of invariants gives a degreewise surjection from $B$ to the
enhanced permanent algebra. In either case, part~1 gives
\begin{equation}\label{G26:degreewise_bound}
             m_{j\nu}(\ell^a\operatorname{per}_m)\le\dim B_j.
\end{equation}

Both $A$ and $B$ have a stable point with the no-scaling property.
For $A$, apply R-28 with degree $n$, matrix size $n-1$, row
count $r$, and the direction $\nu$: its strict condition is
$\nu_1=as>s$. The hypotheses $r\le n-1$ hold since
$r\le m<n$. For $B$, use degree $m$, matrix size $m-1$,
row count $q=r-1$, and the direction $\mu$. Its strict
condition is $\mu_1>s$, and $q\le m-1$.
The smaller padded permanent embeds in $\mathcal P_{m,q}$
by block diagonal padding with one linear entry. At $q=3$
the same point also lies in the full ternary coefficient space.
The common-point clause of R-28 applies in row count three
and does not require size at least seventeen for that clause.
The enhanced representations, and their scalar characters,
are exactly those just specified.

It follows from G-22's extraction and the stable quotient
dimension formula that
\[
                        \dim A=u,\qquad\dim B=v.
\]
Here R-28 supplies the determinant image dimension, G-21
supplies the permanent dimension for $q\ge4$, and the ternary
ambient dimension is $\binom{m+2}{2}$ for $q=3$.
No common invariant between two different groups is being
asserted. The no-scaling property gives gcd one for the
nonzero degrees in each ring separately, by the closed-orbit
separation argument already proved in R-28.

\paragraph{3. The determinant quotient has strictly larger dimension.}
Subtracting a constant from every part does not change flag
dimension. Deleting the first part removes exactly its strict
comparisons with the others, so
\[
 r_\eta-r_\theta
       =\#\{j>1:\nu_1>\nu_j\}\ge0.
\]
For $r\ge5$ this gives
\[
 \begin{aligned}
 u-v&=(r-2)n^2-(r-1)m^2+2m-2r+1+r_\eta-r_\theta\\
 &\ge\frac{r^2-3r+1}{(r-1)^2}m^2+2m-2r+1>0,
 \end{aligned}
\]
since $m\ge r$ and $r^2-3r+1>0$. For $r=4$,
\[
 u-v=2n^2-\binom{m+2}{2}-5+r_\eta-r_\theta
       \ge\frac{55m^2-27m-108}{18}>0.
\]
The explicit formulas give $v\ge2$; for instance at $r\ge5$
one has $v\ge3m^2-2m+2r+2>2$ by $r\le m$, and the ternary
case is immediate. Thus $u\ge v+1$.

\paragraph{4. Effective comparison in every original stretch.}
The two ambient representation dimensions are $b_D,b_B$.
Their action degrees are at most $L$: the smaller one has
degree $ms+(q-1)e'$, and
$e-e'=\nu_1-\nu_r\ge0$. Derksen's bounds used in G-23
therefore give the displayed common bound $\beta$ for the
two invariant generator degrees, even though the groups and
representations differ. Their numbers of generators are each
at most $T$. The identical generator grouping proves that
their $E$th Veroneses are standard graded. Each no-scaling
point gives nonzero elements in every degree at least
$\beta^2$, and $E\ge\beta^2$.

The determinant enhanced cone has its degree-$p$ dense
homogeneous parameter map from G-23. For $r\ge5$, the lower
permanent map has parameter degree $ms+e'\le p$. For $r=4$,
the full coefficient-space parameter has degree $s+e'\le p$.
The base-locus-aware image-degree bound thus gives
\[
                  \deg\operatorname{Proj} B^{(E)}\le(pE)^{v-1}.
\]
Write $j=Ek+t$, $0\le t<E$. In each ring separately, multiply
by its nonzero elements of degrees $E+t$ and $2E-t$ to obtain
the same two-sided Hilbert sandwich as G-23. Linear Noether
normalization and the Chardin upper bound yield
\[
 \dim A_j\ge\binom{k+v-1}{v},\qquad
 \dim B_j\le(pE)^{v-1}\binom{k+v+1}{v-1}.
\]
Their binomial ratio is at least $k/(4v)$ for $k\ge v+1$.
Condition \eqref{G26:threshold} implies
$k>4v(pE)^{v-1}$, so the first bound is strictly larger.
Together with \eqref{G26:degreewise_bound}, this proves the
equality-face conclusion at every stated integer stretch.

If instead $\nu_1<as$, part~1 makes every padded multiplicity
zero. The determinant multiplicity is eventually positive:
$\nu_1\ge ns/r>s$ since $r\le m<n$, so the first R-28
application in part~2 still applies. Its generator-semigroup
bound gives occurrence for every sufficiently large $j$.
Finally combine these two cases with G-23 in the strict
half-space and its proof that $\Delta$ never vanishes. This
proves the exact eventual classification in the theorem.
\end{proof}

\paragraph{Scope.}
The Pieri upper bound is an all-degree reduction and holds
beyond the row range of the eventual comparison. The strict
boundary comparison proved here uses $4\le r\le m$ and
$m\ge17$. It does not assert a positive-degree lifting from
the lower permanent algebra, eliminate earlier boundary gaps
when $r\ge5$, or settle boundary faces with more than $m$
rows. It is independent of the conditional analytic inputs.
